\chapter{Research Questions}

The overall objective of this research is to develop and evaluate a workflow from airborne RGB--LiDAR sensing to helicopter-oriented conflict assessment and avoidance. The central research question is:

\begin{quote}
\textbf{How can airborne perception provide sufficiently accurate, robust information about non-cooperative drones to support helicopter detect-and-avoid decisions?}
\end{quote}

This overarching question is divided into three research questions corresponding to the identified research gaps:

\begin{enumerate}
    \item[\textbf{RQ1}] \textbf{How should a synchronized air-to-air RGB--LiDAR dataset and benchmark be designed to enable reproducible evaluation of small-drone detection, localization, and tracking under DAA-relevant conditions?}

    \textit{This question addresses the lack of publicly available airborne LiDAR drone recordings.}

    \item[\textbf{RQ2}] \textbf{How can existing LiDAR detectors be evaluated and improved through a novel architecture for sparse aerial targets?}

    \textit{This question examines whether conventional point-, voxel-, and pillar-based detectors can be transferred to aerial environments and a new improved model architecture can be implemented to improve accuracy and inference time.}

    \item[\textbf{RQ3}] \textbf{How do the uncertainty of the perception system affect conflict assessment and avoidance planning in helicopter--drone encounters?}

    \textit{This question connects the outputs of detection and tracking to the downstream DAA process. It will investigate how predictions of intruder pose and trajectory influence risk estimation and the generation of avoidance maneuvers.}
\end{enumerate}

Together, the three research questions define a sequential structure. The first establishes the required data and evaluation basis, the second investigates suitable methods for estimating the intruder state, and the third evaluates how the resulting perception performance affects safety-relevant helicopter DAA.

